% Options for packages loaded elsewhere
\PassOptionsToPackage{unicode}{hyperref}
\PassOptionsToPackage{hyphens}{url}
\PassOptionsToPackage{dvipsnames,svgnames,x11names}{xcolor}
%
\documentclass[
  letterpaper,
  DIV=11,
  numbers=noendperiod,
  oneside]{scrreprt}

\usepackage{amsmath,amssymb}
\usepackage{iftex}
\ifPDFTeX
  \usepackage[T1]{fontenc}
  \usepackage[utf8]{inputenc}
  \usepackage{textcomp} % provide euro and other symbols
\else % if luatex or xetex
  \usepackage{unicode-math}
  \defaultfontfeatures{Scale=MatchLowercase}
  \defaultfontfeatures[\rmfamily]{Ligatures=TeX,Scale=1}
\fi
\usepackage{lmodern}
\ifPDFTeX\else  
    % xetex/luatex font selection
\fi
% Use upquote if available, for straight quotes in verbatim environments
\IfFileExists{upquote.sty}{\usepackage{upquote}}{}
\IfFileExists{microtype.sty}{% use microtype if available
  \usepackage[]{microtype}
  \UseMicrotypeSet[protrusion]{basicmath} % disable protrusion for tt fonts
}{}
\makeatletter
\@ifundefined{KOMAClassName}{% if non-KOMA class
  \IfFileExists{parskip.sty}{%
    \usepackage{parskip}
  }{% else
    \setlength{\parindent}{0pt}
    \setlength{\parskip}{6pt plus 2pt minus 1pt}}
}{% if KOMA class
  \KOMAoptions{parskip=half}}
\makeatother
\usepackage{xcolor}
\usepackage[left=1in,marginparwidth=2.0666666666667in,textwidth=4.1333333333333in,marginparsep=0.3in]{geometry}
\setlength{\emergencystretch}{3em} % prevent overfull lines
\setcounter{secnumdepth}{5}
% Make \paragraph and \subparagraph free-standing
\ifx\paragraph\undefined\else
  \let\oldparagraph\paragraph
  \renewcommand{\paragraph}[1]{\oldparagraph{#1}\mbox{}}
\fi
\ifx\subparagraph\undefined\else
  \let\oldsubparagraph\subparagraph
  \renewcommand{\subparagraph}[1]{\oldsubparagraph{#1}\mbox{}}
\fi


\providecommand{\tightlist}{%
  \setlength{\itemsep}{0pt}\setlength{\parskip}{0pt}}\usepackage{longtable,booktabs,array}
\usepackage{calc} % for calculating minipage widths
% Correct order of tables after \paragraph or \subparagraph
\usepackage{etoolbox}
\makeatletter
\patchcmd\longtable{\par}{\if@noskipsec\mbox{}\fi\par}{}{}
\makeatother
% Allow footnotes in longtable head/foot
\IfFileExists{footnotehyper.sty}{\usepackage{footnotehyper}}{\usepackage{footnote}}
\makesavenoteenv{longtable}
\usepackage{graphicx}
\makeatletter
\def\maxwidth{\ifdim\Gin@nat@width>\linewidth\linewidth\else\Gin@nat@width\fi}
\def\maxheight{\ifdim\Gin@nat@height>\textheight\textheight\else\Gin@nat@height\fi}
\makeatother
% Scale images if necessary, so that they will not overflow the page
% margins by default, and it is still possible to overwrite the defaults
% using explicit options in \includegraphics[width, height, ...]{}
\setkeys{Gin}{width=\maxwidth,height=\maxheight,keepaspectratio}
% Set default figure placement to htbp
\makeatletter
\def\fps@figure{htbp}
\makeatother

\KOMAoption{captions}{tableheading}
\makeatletter
\@ifpackageloaded{bookmark}{}{\usepackage{bookmark}}
\makeatother
\makeatletter
\@ifpackageloaded{caption}{}{\usepackage{caption}}
\AtBeginDocument{%
\ifdefined\contentsname
  \renewcommand*\contentsname{Table of contents}
\else
  \newcommand\contentsname{Table of contents}
\fi
\ifdefined\listfigurename
  \renewcommand*\listfigurename{List of Figures}
\else
  \newcommand\listfigurename{List of Figures}
\fi
\ifdefined\listtablename
  \renewcommand*\listtablename{List of Tables}
\else
  \newcommand\listtablename{List of Tables}
\fi
\ifdefined\figurename
  \renewcommand*\figurename{Figure}
\else
  \newcommand\figurename{Figure}
\fi
\ifdefined\tablename
  \renewcommand*\tablename{Table}
\else
  \newcommand\tablename{Table}
\fi
}
\@ifpackageloaded{float}{}{\usepackage{float}}
\floatstyle{ruled}
\@ifundefined{c@chapter}{\newfloat{codelisting}{h}{lop}}{\newfloat{codelisting}{h}{lop}[chapter]}
\floatname{codelisting}{Listing}
\newcommand*\listoflistings{\listof{codelisting}{List of Listings}}
\makeatother
\makeatletter
\makeatother
\makeatletter
\@ifpackageloaded{caption}{}{\usepackage{caption}}
\@ifpackageloaded{subcaption}{}{\usepackage{subcaption}}
\makeatother
\makeatletter
\@ifpackageloaded{sidenotes}{}{\usepackage{sidenotes}}
\@ifpackageloaded{marginnote}{}{\usepackage{marginnote}}
\makeatother
\ifLuaTeX
  \usepackage{selnolig}  % disable illegal ligatures
\fi
\usepackage{bookmark}

\IfFileExists{xurl.sty}{\usepackage{xurl}}{} % add URL line breaks if available
\urlstyle{same} % disable monospaced font for URLs
\hypersetup{
  pdftitle={FT de Portas Abertas 2025},
  pdfauthor={Comissão},
  colorlinks=true,
  linkcolor={blue},
  filecolor={Maroon},
  citecolor={Blue},
  urlcolor={Blue},
  pdfcreator={LaTeX via pandoc}}

\title{FT de Portas Abertas 2025}
\author{Comissão}
\date{2025-01-03}

\begin{document}
\maketitle

\renewcommand*\contentsname{Table of contents}
{
\hypersetup{linkcolor=}
\setcounter{tocdepth}{2}
\tableofcontents
}
\bookmarksetup{startatroot}

\chapter{Introdução}\label{introduuxe7uxe3o}

\url{https://www.youtube.com/watch?v=M8yzSSB4ko4}

\begin{figure*}

FT Portas abertas

\begin{figure}[H]

{\centering \includegraphics{index_files/mediabag/organizacao/organizacao.pdf}

}

\caption{Organização}

\end{figure}%

\begin{figure}[H]

{\centering \includegraphics{index_files/mediabag/justificativa/justificativa.pdf}

}

\caption{Justificativa}

\end{figure}%

\begin{figure}[H]

{\centering \includegraphics{index_files/mediabag/sobre/sobre.pdf}

}

\caption{Sobre}

\end{figure}%

\begin{figure}[H]

{\centering \includegraphics{index_files/mediabag/localizacao/localizacao.pdf}

}

\caption{Como chegar}

\end{figure}%

\begin{figure}[H]

{\centering \includegraphics{index_files/mediabag/atividades/atividades.pdf}

}

\caption{Atividades}

\end{figure}%

\begin{figure}[H]

{\centering \includegraphics{index_files/mediabag/publico-alvo/publico-alvo.pdf}

}

\caption{Público-alvo}

\end{figure}%

\begin{figure}[H]

{\centering \includegraphics{index_files/mediabag/inscricao/inscricao.pdf}

}

\caption{Inscrição}

\end{figure}%

\begin{figure}[H]

{\centering \includegraphics{index_files/mediabag/custos/custos.pdf}

}

\caption{Custos}

\end{figure}%

\end{figure*}%

\bookmarksetup{startatroot}

\chapter{Organização}\label{organizauxe7uxe3o}

Portaria FT/DIR nº 16/2024, de 16/12/2024.

\begin{itemize}
\item
  Prof.~Dr.~Vitor Rafael Coluci (Presidente)
\item
  Profa Dra Marta Siviero Guilherme Pires
\item
  Prof Dr Murilo Cesar Lucas
\item
  Prof Dr Felippe Benavente Canteras
\item
  Prof Dr Yuri Alexandre Meyer
\item
  Prof Dr Ivan de Oliveira
\item
  Profa. Dra. Vanessa Testoni
\item
  Prof Dr Plínio Roberto Souza Vilela
\item
  Prof Dr Ulysses Martins Dias
\end{itemize}

\bookmarksetup{startatroot}

\chapter{Público-alvo}\label{puxfablico-alvo}

A \emph{FT de Portas Abertas} tem como público-alvo principalmente
\textbf{jovens que estão cursando o ensino médio} e que estão na fase da
escolha de um curso universitário.

De acordo com
\href{https://qedu.org.br/municipio/3526902-limeira/censo-escolar}{dados
de 2023} do portal de dados educacionais
\href{https://qedu.org.br/}{QEdu}, Limeira teve mais de \textbf{3000}
estudantes matriculados na 3\textsuperscript{o} ano do Ensino Médio,
sendo 2654 matrículas 3a seria de 25 escolas públicas e 562 matrículas
de 18 escolas particulares. Para esses estudantes, a \emph{FT de Portas
Abertas} apresentará opções de cursos gratuitos, de qualidade e
associados a áreas muito relevantes atualmente e que estão disponíveis
na cidade de Limeira.

Além de estudantes do Ensino Médio, a \emph{FT de Portas Abertas} também
pretende alcançar a \textbf{população de forma geral}, atingindo, por
exemplo, os familiares desses estudantes, a fim de divulgar atividades
de pós-graduação, de pesquisa e de extensão que são oferecidos pela
Faculdade de Tecnologia.

\bookmarksetup{startatroot}

\chapter{Justificativa}\label{justificativa}

A Faculdade de Tecnologia (FT) oferece 6 cursos de graduação sendo 2
deles de Tecnologia, 3 de Engenharia e 1 de Bacharelado. Os cursos são:

\begin{itemize}
\item
  \textbf{Tecnologia}: Tecnologia em Análise e Desenvolvimento de
  Sistemas e Tecnologia em Saneamento Ambiental;
\item
  \textbf{Engenharias}: Engenharia Ambiental, Engenharia de
  Telecomunicações e Engenharia de Transportes;
\item
  \textbf{Bacharelado}: Sistemas de Informação.
\end{itemize}

Esses cursos atendem atualmente

De acordo com
\href{https://www.comvest.unicamp.br/base-de-dados-comvest/vestibulares/vestibulares-anteriores/}{dados}
da Comissão Permanente para os Vestibulares
(\href{https://www.comvest.unicamp.br/}{Convest}) da Unicamp, a procura
pelos cursos oferecidos pela FT tem sido majoritariamente \emph{baixa}.
Essa procura pode ser quantificada pela relação candidato-vaga para a
1\textsuperscript{a} fase do vestibular, indicada nas figuras a seguir.

\begin{figure}[H]

{\centering \includegraphics[width=0.6\textwidth,height=\textheight]{justificativa/../dados/grafico-candidato-vaga-engenharias.png}

}

\caption{Evolução da relação de candidatos por vaga para os cursos de
Engenharia da FT.}

\end{figure}%

\begin{figure}[H]

{\centering \includegraphics[width=0.6\textwidth,height=\textheight]{justificativa/../dados/grafico-candidato-vaga-bacharelado.png}

}

\caption{Evolução da relação de candidatos por vaga para o curso de
Bacharelado da FT.}

\end{figure}%

\begin{figure}[H]

{\centering \includegraphics[width=0.6\textwidth,height=\textheight]{justificativa/../dados/grafico-candidato-vaga-tecnologia.png}

}

\caption{Evolução da relação de candidatos por vaga para os cursos de
Tecnologia da FT.}

\end{figure}%

Os cursos da área de computação (dados em azul) tem apresentado um
crescimento da relação candidato-vaga, principlamente para o curso de
Tecnologia em Análise de Sistemas que alcançou 24 candidatos por vaga em
2024, ficando entre os
\href{https://unicamp.br/noticias/2024/09/20/comvest-registra-63-mil-estudantes-inscritos-no-vestibular-2025/}{10
mais concorridos da Unicamp}.

Por outro lado, esse desempenho não tem ocorrido para os cursos de
Engenharia, os quais tem apresentado, na média, um decréscimo da relação
candidato-vaga. Em 2024, essa relação ficou em torno de 5 candidatos por
vaga para os 3 cursos. Esse valor pode ser comparado com as relações dos
cursos de Engenharia de Produção e de Manufatura, ambos oferecidos
também em Limeira mas no Campus II, de 17 e 4, respectivamente.

Essa comparação indica que aspectos do próprio curso influenciam a
procura, uma vez que cursos oferecidos no mesmo campus possuem relações
(4 e 17) tão diferentes. No entanto, por outro lado, principalmente para
o caso da FT, há uma percepção por parte do corpo docente da FT de que o
desconhecimento pela população de Limeira sobre a FT influencia a
procura por seus cursos. Essa percepção está associada ao Cotil, o
colégio técnico da Unicamp em Limeira, que divide o Campus I com a FT.
Parece haver a impressão que no Campus I só existe o Cotil, uma escola
técnica, desconhecendo-se a FT, uma faculdade de Ensino Superior. É
comum e recorrente a confusão de levar professores externos que viriam
para a FT partipar de bancas, por exemplo, para o Campus II. Isso porque
a população em geral e, particularmente, profissionais como motoristas e
taxistas, associarem o Campus II como ``Unicamp'' e o Campus I apenas
como ``Cotil''.

A FT também participa da \textbf{Unicamp de Portas Abertas}
(\href{https://upa.unicamp.br/}{UPA}) onde os
\href{https://wordpress.ft.unicamp.br/upa/}{cursos de graduação são
apresentados}. Essa participação tem beneficiado a divulgação dos cursos
da FT. No entanto, pelo fato da FT estar localizada numa outra cidade,
apenas a participação do evento em Campinas não atinge um público-alvo
importante para o crescimento da procura pelos cursos da FT que são os
alunos de Ensino Médio de Limeira e região. A falta de contato com as
instalações e atividades realizadas na FT restringe as oportunidades
para que possíveis ingressantes da Unicamp optem pelos cursos da FT.

Pelas razões expostas aqui, estamos propondo a realização da \emph{FT de
Portas Abertas} no próprio campus da FT, a fim expandir o conhecimento
da população de Limeira e região sobre a FT e as oportunidades de
Ensino, Pesquisa e Extensão que ela oferece. Esse tipo de ação já é
feito pelos
\href{https://unicamp.br/noticias/2024/07/24/colegios-tecnicos-da-unicamp-abrem-portas-para-novos-estudantes/}{colégios
técnicos da Unicamp}, que ficam foram do campus de Barão Geraldo.

\bookmarksetup{startatroot}

\chapter{Custos}\label{custos}

Os custos associados ao \emph{FT de Portas} se referem a itens para a
organização do evento (camisetas) e para a divulgação dos cursos na
forma de \emph{folders} a serem entregues para os visitantes.

O uso de um mapa ``lúdico'', similar ao
\href{https://internationaloffice.unicamp.br/wp-content/uploads/sites/26/2023/10/MAPA-BARAO-GERALDO-baixa.jpg\#/home}{apresentado
pela própria Unicamp} em seu
\href{https://unicamp.br/unicamp/mapas/}{sítio}, é uma proposta para
apresentar e enfatizar a FT dentro do Campus I da Unicamp em Limeira e,
assim, separar a associação cultural da população limeirense de que o
campus II (FCA) é a Unicamp e o campus I (FT e Cotil) é o somente o
Cotil.

\begin{longtable}[]{@{}
  >{\raggedright\arraybackslash}p{(\columnwidth - 8\tabcolsep) * \real{0.1212}}
  >{\raggedright\arraybackslash}p{(\columnwidth - 8\tabcolsep) * \real{0.1818}}
  >{\raggedright\arraybackslash}p{(\columnwidth - 8\tabcolsep) * \real{0.3182}}
  >{\raggedright\arraybackslash}p{(\columnwidth - 8\tabcolsep) * \real{0.1818}}
  >{\raggedright\arraybackslash}p{(\columnwidth - 8\tabcolsep) * \real{0.1970}}@{}}
\toprule\noalign{}
\begin{minipage}[b]{\linewidth}\raggedright
Item
\end{minipage} & \begin{minipage}[b]{\linewidth}\raggedright
Descrição
\end{minipage} & \begin{minipage}[b]{\linewidth}\raggedright
Valor unitário (R\$)
\end{minipage} & \begin{minipage}[b]{\linewidth}\raggedright
Quantidade
\end{minipage} & \begin{minipage}[b]{\linewidth}\raggedright
Total (R\$)
\end{minipage} \\
\midrule\noalign{}
\endhead
\bottomrule\noalign{}
\endlastfoot
1 & Camiseta organizador & 100 & 12 & 1200 \\
2 & Camiseta suporte & 80 & 20 & 1600 \\
3 & \href{./display-para-festas-4_1-150x150.jpg}{Display para foto} &
1000 & 2 & 2000 \\
4 &
\href{https://unicamp.br/unicamp/sites/default/files/2023-03/Atu_20220915_Calourada-2023_Scarpa_AJS_8576_Capa.jpg}{Banner
de recepção} & 1000 & 1 & 1000 \\
5 & Confeçcão de \href{mapa-zoo.jpg}{Mapa lúdico} & 8000 & 1 & 8000 \\
6 & Folder: Mapa lúdico (frente) + infos de cursos (trás) & 10? & 500? &
5000 \\
7 & Equipamento A & XXX & 1 & XXX \\
8 & Equipamento B & XXX & 1 & XXX \\
9 & Equipamento C & XXX & 1 & XXX \\
Total Geral & & & & \textbf{18800} \\
\end{longtable}

\bookmarksetup{startatroot}

\chapter{Atividades}\label{atividades}

\bookmarksetup{startatroot}

\chapter{Localização}\label{localizauxe7uxe3o}

\begin{verbatim}
<folium.folium.Map at 0x132b709b0>
\end{verbatim}

Faculdade de Tecnologia/UNICAMP, Rua Paschoal Marmo, 188 Jardim Nova
Itália - Limeira-SP-BR

Coordenadas: -22.562142069, -47.42408844



\end{document}
